\subsection*{Appendix E: How to Interview Someone Who Rejected Your Price}

A strong lost-buyer interview does not begin, ``Was it the price?'' That question invites a polite and often misleading answer. First ask the prospect to reconstruct the decision as if describing an actual event to a colleague. What triggered the purchase? When did they 1st consider a solution? What were the alternatives? What would doing nothing cost them? What changed between the 1st interest and the final decision? Who else was involved? What did they ultimately do?

Only then discuss price. Was the quoted fee larger than the relevant alternative? Was it beyond an allocated budget? Was the total time investment too large? Was the downside of choosing an unproven provider unacceptable? Would a narrower scope, different evidence, or different payer have changed the choice? The interviewer should separate explanations volunteered before probing from explanations produced after the founder suggests them.

The useful output is a purchasing chronology, not a binary ``too expensive'' label. It should show a real next step: improve proof, reduce buyer work, create a complete smaller product, change the target segment, change payer, or accept that the market is adequately served.

\subsection*{Appendix F: Seven Difficult Founder Exercises}

\paragraph*{1. The substitute inversion.} Select a service you believe is valuable. Write down 5 substitutes, including doing nothing. For each, explain the actual purchase decision that could cause a buyer to prefer it. Which competitor would you never mention in a sales pitch because it makes your value harder to justify? Investigate that competitor 1st.

\paragraph*{2. The full-sacrifice audit.} Add nonmonetary adoption burden to the advertised fee: preparation time, switching, tool purchases, social friction, probability of failure, and learning curve. Which benefits still survive? Could you reduce the sacrifice more cheaply than adding value?

\paragraph*{3. The price-reference audit.} Interview 10 recent buyers about what price they expected before seeing your offer and where that number came from. Do not ask whether your price is fair. Determine whether their comparison category is the 1 you believe you occupy.

\paragraph*{4. The payer switch.} Draw 3 versions of the same benefit: individual pays, employer pays, institution pays. For each specify budget, proof, approval, customization, and cash timing. Which option has the most value for the payer, not merely the user?

\paragraph*{5. The credible-premium test.} Remove the founder's biography, all prestige adjectives, and any promises of future network value. What can a skeptical buyer still inspect today? Would you purchase at the fee you propose if you did not know the founder personally?

\paragraph*{6. The unsafe-affordability test.} Identify the buyer most likely to value the product but least able to pay safely. Can you design a valid sponsorship, employer-funded, scholarship, or lower-cost route without transferring unacceptable risk or degrading core quality?

\paragraph*{7. The longitudinal test.} Ask buyers after experience whether they would repeat the same choice. Compare their answer with actual outcomes, complaints, and referrals. If the customer was satisfied but did not achieve the promised capability, which part of your WTP theory was wrong?

\subsection*{Appendix G: A Comparative Evidence Sheet}

Build 1 row for every prospect who encounters a real price. Capture anonymized segment and trigger; alternative; stated range if volunteered; full offer shown; actual fee and terms; whether payment was made; delay to decision; objection; follow-up action; whether the person actually used the product; and postpurchase experience. Preserve the distinction between private customer evidence and permissioned public testimonials.

A field researcher should review the data by \emph{buyer situation} and \emph{offer}, not merely aggregate it into 1 mean. If all paid buyers came from the founder's personal network and all cold prospects refused, the appropriate conclusion is not ``the market pays.'' It is that a particular trust channel produced early customers, and broader demand remains to be tested.

\subsection*{Appendix H: The Economic and Moral Boundary of WTP}

A product can have high willingness to pay because a buyer has almost no alternatives and faces frightening consequences. Emergency services, high-stakes credentials, healthcare, and financial products can exhibit such situations. A founder's ability to charge more is not sufficient evidence that doing so is good strategy or ethical practice. Contracts, professional duties, competition policy, and consumer protection may constrain the transaction.

An enduring institution should ask 3 additional questions: Would the buyer consider the price defensible after becoming fully informed? Would the price allow the firm to maintain or improve quality? Would the arrangement preserve enough buyer welfare and public legitimacy to sustain relationships over years? The long game is played not against the customer but with the customer's repeated consent.
